% Input after extension.tex; requires amsmath and amsthm.
\subsection{一维无限质量权的内禀临界发散}\label{subsec:one-dimensional-critical}

下述结论只处理等指数乘积；它给出内禀模型在一维无限质量权下的临界区间，
并不推广到高维或有限总质量。

\begin{theorem}\label{thm:one-dimensional-infinite-critical}
令 $m\ge2$，$0\le w\in L^1_{\rm loc}(\mathbb R)$ 非零，
$\mu=w\,dx$ 且 $\mu(\mathbb R)=\infty$。对
\[
0<\lambda\le\lambda_m^*:=\left(1-\frac1m\right)^m,
\]
共同中心、公共交集归一化的内禀多线性 Tauberian 常数满足
\[
C^{(m),\rm int}_{c,1,w}(\lambda)=\infty.
\]
\end{theorem}

\begin{proof}
令 $T(x)=\int_0^xw(t)\,dt$。这是连续非减函数。令
$J=(-a,b)$ 为 $T(\mathbb R)$ 的内部（端点是否属于像集无关紧要），则
$a+b=\mu(\mathbb R)=\infty$。对 $u<v$ 属于 $J$，由连续性取
$y<z$ 使 $T(y)=u,T(z)=v$；$T^{-1}((u,v))$ 与 $(y,z)$ 至多相差
若干 $T$-纤维，而每个这样的纤维的 $\mu$-质量为零。因此
\[
 \mu\bigl(T^{-1}((u,v))\bigr)=\mu((y,z))=v-u.
\tag{C.1}
\]
换言之，$T_\#\mu$ 是 $J$ 上的 Lebesgue 测度。这个表述已包含
$w=0$ 所产生的平坦区间：它们是 $T$ 的纤维，故有零 $\mu$-质量。

在 $J$ 上按周期 $1$ 划分
\[
\widetilde A_i=\bigcup_{k\in\mathbb Z}\left[k+\frac{i-1}{m},
k+\frac{i}{m}\right)\cap J,\qquad 1\le i\le m,
\]
并令 $A_i=T^{-1}(\widetilde A_i)$；若 $T$ 取到 $J$ 的端点，则将相应
零质量纤维任意指定给某一个 $A_i$，从而 $\{A_i\}_{i=1}^m$ 是严格的分割。
对任意物理区间 $Q$，端点截断给出
\[
\left|\mu(A_i\cap Q)-\frac{\mu(Q)}m\right|\le2.\tag{C.2}
\]
确实，按 \textup{(C.1)}，左端是 $\widetilde A_i$ 在一个 Lebesgue
区间上的计数误差；去掉两端至多两个不完整的长度 $1$ 周期后，余下完整
周期的误差为零，而所去部分的绝对误差至多为 $2$。
事实上，$J$ 至少有一个无穷端，而每个 $\widetilde A_i$ 在该端都占
无限 Lebesgue 测度，故每个 $A_i$ 都有无限 $\mu$-质量。由
$A_1\cap[-L,L]\uparrow A_1$，可取 $L<\infty$ 使
$F:=A_1\cap[-L,L]$ 满足 $0<f:=\mu(F)<\infty$。设
$E_i=A_i^c\cup F$。则
$\cap_iE_i=F$。

固定 $x\in\mathbb R$，令 $Q_R=(x-R,x+R)$。函数
$M_R:=\mu(Q_R)=T(x+R)-T(x-R)$ 连续、非减且趋于无穷。先取 $R_0$
使 $F\subset Q_{R_0}$，再任选整数 $N>M_{R_0}$；介值性给出某个
$R\ge R_0$ 使 $M_R=N$。由于 $\widetilde A_i$ 的指示函数以 $1$ 为周期、
一个周期内积分为 $1/m$，其在任意长度为整数 $N$ 的区间上的积分正好为
$N/m$。结合 \textup{(C.1)}，
\[
 \mu(A_i\cap Q_R)=\frac Nm\qquad(1\le i\le m).
\tag{C.3}
\]
注意 $F\subset A_1$，故它只给第一个补集带来新增质量：
\[
 \frac{\mu(E_1\cap Q_R)}N=c+\frac fN,
 \qquad
 \frac{\mu(E_i\cap Q_R)}N=c\quad(2\le i\le m),
 \qquad c=1-\frac1m.
\]
于是该共同中心区间严格见证
\[
 \prod_{i=1}^m\frac{\mu(E_i\cap Q_R)}{\mu(Q_R)}
=c^{m-1}\left(c+\frac fN\right)>c^m.
\tag{C.4}
\]
故每个 $x$ 属于
$\{\mathcal M^{(m)}_{c,w}(\mathbf1_{E_1},\ldots,\mathbf1_{E_m})>\lambda\}$。
该光晕为全空间，且其 $\mu$-质量无穷；除以 $\mu(F)=f$ 即得结论。
\end{proof}

\begin{remark}
本证明的严格余量来自公共集合 $F$ 的总质量，而不是周期条带的极限等号。
若 $\mu(\mathbb R)<\infty$，最后一步不能由光晕的无穷质量推出；该情形须另行处理。
\end{remark}

\begin{corollary}
一维时，对任意非零非负局部可积权和 $m\ge2$，在 $0<\lambda\le\lambda_m^*$ 内，内禀共同中心多集合常数均为无穷。
\end{corollary}
\begin{proof}无限总质量用本节定理，有限总质量用命题 \ref{prop:finite-mass-critical} 的等指数特例。\end{proof}
